\documentclass[10pt,a4paper]{amsart}
\usepackage{amsmath,amssymb,amsthm,graphicx}
\usepackage{verbatim}
\newenvironment{reference}{\comment}{\endcomment}
%\newenvironment{reference}{}{}
\newenvironment{slogan}{\comment}{\endcomment}
\newenvironment{history}{\comment}{\endcomment}

% For commutative diagrams we use Xy-pic
\usepackage[all]{xy}

% We use 2cell for 2-commutative diagrams.
\xyoption{2cell}
% Theorem environments.
%
\theoremstyle{plain}
\newtheorem{theorem}[subsection]{Theorem}
\newtheorem{proposition}[subsection]{Proposition}
\newtheorem{lemma}[subsection]{Lemma}

\theoremstyle{definition}
\newtheorem{definition}[subsection]{Definition}
\newtheorem{example}[subsection]{Example}
\newtheorem{exercise}[subsection]{Exercise}
\newtheorem{situation}[subsection]{Situation}

\theoremstyle{remark}
\newtheorem{remark}[subsection]{Remark}
\newtheorem{remarks}[subsection]{Remarks}

\numberwithin{equation}{subsection}

% Macros
%
\def\lim{\mathop{\mathrm{lim}}\nolimits}
\def\colim{\mathop{\mathrm{colim}}\nolimits}
\def\Spec{\mathop{\mathrm{Spec}}}
\def\Hom{\mathop{\mathrm{Hom}}\nolimits}
\def\Ext{\mathop{\mathrm{Ext}}\nolimits}
\def\SheafHom{\mathop{\mathcal{H}\!\mathit{om}}\nolimits}
\def\SheafExt{\mathop{\mathcal{E}\!\mathit{xt}}\nolimits}
\def\Sch{\mathit{Sch}}
\def\Mor{\mathop{\mathrm{Mor}}\nolimits}
\def\Ob{\mathop{\mathrm{Ob}}\nolimits}
\def\Sh{\mathop{\mathit{Sh}}\nolimits}
\def\NL{\mathop{N\!L}\nolimits}
\def\CH{\mathop{\mathrm{CH}}\nolimits}
\def\proetale{{pro\text{-}\acute{e}tale}}
\def\etale{{\acute{e}tale}}
\def\QCoh{\mathit{QCoh}}
\def\Ker{\mathop{\mathrm{Ker}}}
\def\Im{\mathop{\mathrm{Im}}}
\def\Coker{\mathop{\mathrm{Coker}}}
\def\Coim{\mathop{\mathrm{Coim}}}

% Boxtimes
%
\DeclareMathSymbol{\boxtimes}{\mathbin}{AMSa}{"02}

%
% Macros for moduli stacks/spaces
%
\def\QCohstack{\mathcal{QC}\!\mathit{oh}}
\def\Cohstack{\mathcal{C}\!\mathit{oh}}
\def\Spacesstack{\mathcal{S}\!\mathit{paces}}
\def\Quotfunctor{\mathrm{Quot}}
\def\Hilbfunctor{\mathrm{Hilb}}
\def\Curvesstack{\mathcal{C}\!\mathit{urves}}
\def\Polarizedstack{\mathcal{P}\!\mathit{olarized}}
\def\Complexesstack{\mathcal{C}\!\mathit{omplexes}}
% \Pic is the operator that assigns to X its picard group, usage \Pic(X)
% \Picardstack_{X/B} denotes the Picard stack of X over B
% \Picardfunctor_{X/B} denotes the Picard functor of X over B
\def\Pic{\mathop{\mathrm{Pic}}\nolimits}
\def\Picardstack{\mathcal{P}\!\mathit{ic}}
\def\Picardfunctor{\mathrm{Pic}}
\def\Deformationcategory{\mathcal{D}\!\mathit{ef}}
\usepackage{cleveref}
\crefname{originalreference}{original reference}{original references}
\Crefname{originalreference}{Original reference}{Original references}
\makeatletter
\newlabel{example-decompose-P1}{{0FG6}{source}{}{sourcepointer-1}{}}
\newlabel{example-decompose-P1@cref}{{[originalreference][0][]0FG6}{[1][0][]source}{}{}{}}
\makeatother
\crefname{originalreference}{original reference}{original references}
\Crefname{originalreference}{Original reference}{Original references}
\makeatletter
\newlabel{lemma-contravariant-functor}{{0FG3}{source}{}{sourcepointer-2}{}}
\newlabel{lemma-contravariant-functor@cref}{{[originalreference][0][]0FG3}{[1][0][]source}{}{}{}}
\makeatother
\title{The motives of a curve with the zero cycles and with the degree two cycles}
\author{The Stacks Project Authors}
\date{stacks-project:weil.tex:local-snapshot-20261012; GNU FDL 1.2 or later}
\begin{document}
\maketitle
\noindent\textit{Editorial note.} All delivered mathematical text is the author's own, copied from the pinned Stacks Project chapter file weil.tex; no statement or proof step was rewritten, reordered or removed.\par\medskip
\noindent\textit{Editorial note.} The counted proof is the author's complete argument: the construction of the correspondences from the structure morphism and the rational point, and the intersection product computation identifying the compositions with the given idempotents.\par\medskip
\noindent\textit{Editorial note.} The label lemma-inverse-h2 names the counted statement, and the proof environment immediately following it in the source is delivered verbatim as its proof.\par\medskip
\noindent\textit{Editorial note.} The counted statement is carried by the author's own lemma environment in the chapter source.\par\medskip
\noindent\textit{Editorial note.} Pre-open gate: python operations/agent-authoring/tools/precard\_gate.py --source weil.tex --label lemma-inverse-h2 --start 965 --end 1028 --local-id 28133 returned RESULT=OK (L28133 free, occupancy FREE, public interval clear). Pre-assembly tag self-check: canonical tag resolved and not present in reference/public-items.jsonl.\par\medskip
\noindent\textit{Editorial note.} The statement and proof are copied from the pinned chapter file; the cited inputs are kept as external pointers with their canonical Stacks tags.\par\medskip
\noindent\textit{Allowed standard prerequisites (editorial).}\begin{itemize}
\item The functoriality of correspondences and the description of morphisms of motives. Same chapter, the lemma on the contravariant functor, and the description of morphisms between motives as compositions with the involved correspondences. Provides the description of morphisms between the motives under consideration and the identities used to show that the constructed correspondences are mutually inverse.
\item Correspondences and intersection products on products of the curve with itself. The identification of the groups of correspondences of degree zero and one with the groups of cycles and the computation of the intersection product on the triple product. Identifies the compositions of the constructed correspondences with the given idempotents through an explicit intersection product computation.
\end{itemize}\medskip
\ifdefined\hypertarget\hypertarget{sourcepointer-1}{}\fi
\noindent\textit{Source reference.} Original reference 0FG6: sources/stacks-weil/pinned-original-source.tex, source lines 718--718. Prerequisite pointer invoked at proof line(s) 967; it names a standard fact of the same chapter that the counted argument uses, and every step of the counted proof built on it is retained verbatim in the delivered body.\par\medskip
\ifdefined\hypertarget\hypertarget{sourcepointer-2}{}\fi
\noindent\textit{Source reference.} Original reference 0FG3: sources/stacks-weil/pinned-original-source.tex, source lines 610--610. Prerequisite pointer invoked at proof line(s) 979; it names a standard fact of the same chapter that the counted argument uses, and every step of the counted proof built on it is retained verbatim in the delivered body.\par\medskip
\begin{lemma}
\label{lemma-inverse-h2}
With notation as in Example \ref{example-decompose-P1}
\begin{enumerate}
\item
the motive $(X, c_0, 0)$ is isomorphic to the motive
$\mathbf{1} = (\Spec(k), 1, 0)$.
\item
the motive $(X, c_2, 0)$ is isomorphic to the motive
$\mathbf{1}(-1) = (\Spec(k), 1, -1)$.
\end{enumerate}
\end{lemma}

\begin{proof}
We will use Lemma \ref{lemma-contravariant-functor} without further mention.
The structure morphism $X \to \Spec(k)$ gives a correspondence
$a \in \text{Corr}^0(\Spec(k), X)$. On the other hand, the rational
point $x$ is a morphism $\Spec(k) \to X$ which gives a correspondence
$b \in \text{Corr}^0(X, \Spec(k))$. We have $b \circ a = 1$ as a
correspondence on $\Spec(k)$. The composition $a \circ b$ corresponds
to the graph of the composition $X \to x \to X$ which is
$c_0 = [x \times X]$. Thus $a = a \circ b \circ a = c_0 \circ a$
and $b = b \circ a \circ b = b \circ c_0$.
Hence, unwinding the definitions, we see that
$a$ and $b$ are mutually inverse morphisms
$a : (\Spec(k), 1, 0) \to (X, c_0, 0)$ and
$b : (X, c_0, 0) \to (\Spec(k), 1, 0)$.

\medskip\noindent
We will proceed exactly as above to prove the second statement.
Denote
$$
a' \in \text{Corr}^1(\Spec(k), X) = \CH^1(X)
$$
the class of the point $x$. Denote
$$
b' \in \text{Corr}^{-1}(X, \Spec(k)) = \CH_1(X)
$$
the class of $[X]$. We have $b' \circ a' = 1$ as a correspondence
on $\Spec(k)$ because $[x] \cdot [X] = [x]$ on
$X = \Spec(k) \times X \times \Spec(k)$. Computing the
intersection product $\text{pr}_{12}^*b' \cdot \text{pr}_{23}^*a'$
on $X \times \Spec(k) \times X$ gives the cycle
$X \times \Spec(k) \times x$. Hence
the composition $a' \circ b'$ is equal to $c_2$ as a
correspondence on $X$. Thus $a' = a' \circ b \circ a' = c_2 \circ a'$
and $b' = b' \circ a' \circ b' = b' \circ c_2$. Recall that
$$
\Mor((\Spec(k), 1, -1), (X, c_2, 0)) =
c_2 \circ \text{Corr}^1(\Spec(k), X)
\subset
\text{Corr}^1(\Spec(k), X)
$$
and
$$
\Mor((X, c_2, 0), (\Spec(k), 1, -1)) =
\text{Corr}^{-1}(X, \Spec(k)) \circ c_2
\subset
\text{Corr}^{-1}(X, \Spec(k))
$$
Hence, we see that $a'$ and $b'$ are mutually inverse morphisms
$a' : (\Spec(k), 1, -1) \to (X, c_2, 0)$ and
$b' : (X, c_2, 0) \to (\Spec(k), 1, -1)$.
\end{proof}

\end{document}
